%!TEX root = ../main.tex

% Optional: Uncomment if you also have a Chinese abstract
% \begin{abstract}
%   [Your Chinese abstract here]
%   \keywords{关键词1；关键词2；关键词3}
% \end{abstract}

\begin{enabstract}
  Privacy-preserving tabular data synthesis remains challenging due to fundamental trade-offs between data utility, fairness, and formal privacy guarantees. Generative Adversarial Networks (GANs) have demonstrated strong performance in synthesizing mixed-type tabular data but often suffer from training instability, mode collapse, and vulnerability to membership inference attacks. Diffusion models offer improved stability and diversity but face significant challenges when applied directly to tabular data containing discrete and categorical attributes, and they often introduce substantial computational overhead. Furthermore, training-time differential privacy mechanisms such as DP-SGD provide strong theoretical guarantees but frequently degrade downstream utility, particularly for rare categories and minority subgroups.

  This thesis presents \textit{Diff-GAN}, a privacy-preserving tabular data synthesis framework that decouples fidelity-oriented training from privacy enforcement. Diff-GAN employs a lightweight latent-space diffusion refiner that stabilizes GAN training and improves mode coverage without requiring discrete diffusion in data space. Differential privacy is enforced only at model selection time, using a Laplace-noised divergence score with rare-category upweighting to privately select the generator released for synthesis. Experiments on real-world tabular benchmarks, including Adult, Covertype, and an EHR-style dataset, demonstrate that Diff-GAN achieves strong predictive utility, near-zero fairness disparity, improved distribution coverage, and reduced marginal distribution shift compared with strong baseline methods. These results show that Diff-GAN offers a practical and principled solution for high-utility, differentially private tabular data synthesis.

  \enkeywords{Synthetic Data, Tabular Data, GAN, Diffusion Models, Differential Privacy, Private Model Selection, Fairness}
\end{enabstract}